\documentclass[11pt]{article}
\usepackage[a4paper,margin=28mm]{geometry}
\usepackage{amsmath,amssymb,amsthm,mathtools}
\usepackage{booktabs,longtable,array,tabularx}
\usepackage[hidelinks]{hyperref}
\usepackage{microtype}
\usepackage{enumitem}

\newtheorem{theorem}{Theorem}[section]
\newtheorem{proposition}[theorem]{Proposition}
\newtheorem{lemma}[theorem]{Lemma}
\newtheorem{corollary}[theorem]{Corollary}
\theoremstyle{remark}
\newtheorem{remark}[theorem]{Remark}

\newcommand{\Q}{\mathbb Q}
\newcommand{\Z}{\mathbb Z}
\newcommand{\Gzero}{\Gamma_0}
\newcommand{\Gone}{\Gamma_1}
\newcommand{\tht}{\theta_t}
\newcommand{\thq}{\theta_q}
\newcommand{\etaD}[1]{\eta_{#1}}
\newcommand{\chim}{\chi_{-3}}
\newcommand{\chifour}{\chi_{-4}}
\newcommand{\chifive}{\chi_{5}}
\newcommand{\sigin}[2]{\sigma_{#1}^{(#2)}}
\newcommand{\sigout}[2]{\widetilde\sigma_{#1}^{(#2)}}

\title{Eisenstein Sources for the Twelve Modular Sporadic Ap\'ery Companions}
\author{[Author]}
\date{August 2026}

\begin{document}
\maketitle

\begin{abstract}
For the twelve sporadic Ap\'ery-like recurrences admitting modular
parametrizations, we prove that the canonical source of the normalized second
solution is a pure Eisenstein series.  The key simplification is the identity
\[
   \Phi(\tau)=F(\tau)\,\theta_q t(\tau),
\]
where $t$ is the canonical nome coordinate and $F$ is the analytic solution
pulled back to the modular parameter.  Eleven cases are proved uniformly.
Their candidate $t$ and $F$ are eta quotients; Ligozat's criterion proves
modularity, and two small holomorphic modular identities --- a normalized
Wronskian identity and a Rankin--Cohen form of the differential equation ---
prove that these eta quotients are the \emph{canonical} recurrence
parametrizations.  Sturm's theorem then identifies $F\theta_qt$ with the
stated Eisenstein series.  The remaining Ap\'ery $\zeta(2)$ family is treated
on $\Gamma_1(5)$ using Zagier's Rogers--Ramanujan parametrization and a
weight-six square identity.  Thus the previously finite coefficient
identifications become identities for all Fourier coefficients, and no
cuspidal component occurs in any of the twelve sources.
\end{abstract}

\section{The recurrences and the source}

Put $\tht=t\,d/dt$ and $\thq=q\,d/dq$.  We use the two operator families
\begin{align}
L_2&=\tht^2-t(a\tht^2+a\tht+b)+ct^2(\tht+1)^2,
\label{eq:L2}\\
L_3&=\tht^3-t(2\tht+1)(a\tht^2+a\tht+b)+ct^2(\tht+1)^3.
\label{eq:L3}
\end{align}
Their leading polynomials are
\[
P_2(t)=1-at+ct^2,\qquad P_3(t)=1-2at+ct^2.
\]
Let $y_0(t)=1+O(t)$ be the analytic solution, let $q$ be the canonical
Frobenius nome, and write
\[
  t=t(q),\qquad F(q)=y_0(t(q)),\qquad
  K(q):=\thq\log t=\frac{\thq t}{t}.
\]
Thus $t=q+O(q^2)$ and $K=1+O(q)$.

The normalized companion is the unique log-free series $y_B(t)$ with
$y_B(0)=0$ and $L_r y_B=t$.  Exact rectification by the canonical nome gives
\[
 L_r=C(t)F\thq^rF^{-1}.
\]
Since $\thq=K\tht$, comparison of the coefficient of $\tht^r$ gives
\begin{equation}
 C(t)=\frac{P_r(t)}{K^r}.
 \label{eq:C}
\end{equation}
Consequently
\begin{equation}
 y_B=F\thq^{-r}\Phi,
 \qquad
 \Phi=\frac{tK^r}{P_rF}.
 \label{eq:source-general}
\end{equation}
For $r=2$ rectification is automatic from the Frobenius basis.  For $r=3$
we use the symmetric-square identity of Almkvist--Zudilin \cite[Prop.~9]{AZ}:
\begin{equation}
 \widetilde L
 =\tht^2-t\left(2a\tht^2+a\tht+\frac b2\right)
       +ct^2\left(\tht+\frac12\right)^2,
 \label{eq:root}
\end{equation}
whose symmetric square is $L_3$.  If $u_0,u_1$ are its Frobenius solutions,
then $F=u_0^2$ and $u_1/u_0=\log q$.

\begin{proposition}[collapse of the source]\label{prop:collapse}
For both operator families, in the canonical coordinate,
\begin{equation}
 \boxed{\ \Phi=F\,\thq t.\ }
 \label{eq:collapse}
\end{equation}
\end{proposition}

\begin{proof}
For $L_2$, the Wronskian of $F$ and $F\log q$ is
\[
 F^2\frac{d\log q}{dt}=\frac1{tP_2(t)},
\]
with normalization fixed at $t=0$.  Hence $K=P_2F^2$.  Substitution into
\eqref{eq:source-general} gives
\[
 \Phi=\frac{t(P_2F^2)^2}{P_2F}=tP_2F^3
      =F\,\thq t.
\]
For $L_3$, the Wronskian of the two solutions of \eqref{eq:root} is
$1/(t\sqrt{P_3})$.  Since $F=u_0^2$, this gives
$K=F\sqrt{P_3}$, equivalently $K^2=P_3F^2$.  Therefore
\[
 \Phi=\frac{tK^3}{P_3F}=tKF=F\,\thq t.
\]
\end{proof}

The point of \eqref{eq:collapse} is that the source is no longer an opaque
quotient: $\thq t$ is a meromorphic modular form of weight $2$, so a modular
$F$ of weight $r-1$ makes $\Phi$ a form of weight $r+1$.

\section{A uniform finite certificate for eleven eta families}

Write $\eta_d=\eta(d\tau)$, with $q=e^{2\pi i\tau}$.  For an eta quotient
$G=\prod_{d\mid N}\eta_d^{e_d}$ we use the standard Ligozat cusp-order
formula
\begin{equation}
 \operatorname{ord}_{1/c}G=
 \frac{N}{24c\gcd(c,N/c)}
 \sum_{d\mid N}\frac{\gcd(c,d)^2}{d}e_d,
 \qquad c\mid N.
 \label{eq:ligozatorder}
\end{equation}
The usual two congruences
$\sum d e_d\equiv\sum (N/d)e_d\equiv0\pmod{24}$ and
\eqref{eq:ligozatorder} certify the forms below.  Every displayed $t$ has
weight zero, trivial character, and Fourier expansion $q+O(q^2)$; every $F$
is holomorphic, has constant term $1$, and has the character stated below.

\begin{table}[ht]
\centering
\small
\renewcommand{\arraystretch}{1.22}
\begin{tabular}{@{}c c c p{5.1cm} p{5.1cm}@{}}
\toprule
family & $N$ & $\chi_F$ & $t$ & $F$\\
\midrule
$\mathbf A$&6&$\chim$&
$\eta_1^3\eta_6^9/(\eta_2^3\eta_3^9)$&
$\eta_2\eta_3^6/(\eta_1^2\eta_6^3)$\\
$\mathbf B$&36&$\chim$&
$\eta_1^3\eta_4^3\eta_{18}^9/(\eta_2^9\eta_9^3\eta_{36}^3)$&
$\eta_2^9\eta_3\eta_{12}/(\eta_1^3\eta_4^3\eta_6^3)$\\
$\mathbf C$&6&$\chim$&
$\eta_1^4\eta_6^8/(\eta_2^8\eta_3^4)$&
$\eta_2^6\eta_3/(\eta_1^3\eta_6^2)$\\
$\mathbf E$&8&$\chifour$&
$\eta_1^4\eta_4^2\eta_8^4/\eta_2^{10}$&
$\eta_2^{10}/(\eta_1^4\eta_4^4)$\\
$\mathbf F$&12&$\chim$&
$\eta_1^5\eta_3\eta_4^5\eta_6^2\eta_{12}/\eta_2^{14}$&
$\eta_2^{15}\eta_3^2\eta_{12}^2/(\eta_1^6\eta_4^6\eta_6^5)$\\[2pt]
$\alpha$&12&$1$&
$\eta_1^6\eta_3^6\eta_4^6\eta_{12}^6/(\eta_2^{12}\eta_6^{12})$&
$\eta_2^{10}\eta_6^{10}/(\eta_1^4\eta_3^4\eta_4^4\eta_{12}^4)$\\
$\gamma$&6&$1$&
$\eta_1^{12}\eta_6^{12}/(\eta_2^{12}\eta_3^{12})$&
$\eta_2^7\eta_3^7/(\eta_1^5\eta_6^5)$\\
$\delta$&12&$1$&
$\eta_1^4\eta_4^4\eta_6^{16}/(\eta_2^{16}\eta_3^4\eta_{12}^4)$&
$\eta_2^{12}\eta_3\eta_{12}/(\eta_1^3\eta_4^3\eta_6^4)$\\
$\varepsilon$&8&$1$&
$\eta_1^8\eta_8^8/(\eta_2^8\eta_4^8)$&
$\eta_2^6\eta_4^6/(\eta_1^4\eta_8^4)$\\
$\zeta$&9&$1$&
$\eta_1^6\eta_9^6/\eta_3^{12}$&
$\eta_3^{10}/(\eta_1^3\eta_9^3)$\\
$\eta$&20&$\chifive$&
$\eta_1^6\eta_4^6\eta_{10}^{18}/(\eta_2^{18}\eta_5^6\eta_{20}^6)$&
$\eta_2^{15}\eta_5\eta_{20}/(\eta_1^5\eta_4^5\eta_{10}^3)$\\
\bottomrule
\end{tabular}
\caption{The eleven ordinary eta parametrizations.  The first five rows have
$\operatorname{wt}F=1$; the last six have $\operatorname{wt}F=2$.}
\label{tab:eta}
\end{table}

It remains important to prove that Table~\ref{tab:eta} gives the
\emph{canonical} parametrizations of the recurrences, rather than modular
forms whose first few coefficients happen to agree.  This is done by two
small modular identities.

Let $D=\thq$ and $K=D\log t$.  For an eta quotient,
\begin{equation}
 K=\frac1{24}\sum_{d\mid N} e_d d\,E_2(d\tau),
 \label{eq:Keta}
\end{equation}
and, because $t$ is a weight-zero modular function, $K$ is a genuine
holomorphic weight-two modular form.  Define
\begin{align}
H_2&=F D^2F-2(DF)^2+t(-b+ct)P_2(t)F^6,
\label{eq:H2}\\
H_3&=2F D^2F-3(DF)^2+t(-2b+ct)F^4.
\label{eq:H3}
\end{align}
The derivative combinations are Rankin--Cohen brackets, hence modular and
holomorphic.  Formula \eqref{eq:ligozatorder} shows that at every pole of $t$
we have
\begin{equation}
 \operatorname{ord}F=-\operatorname{ord}t.
 \label{eq:polecancel}
\end{equation}
Thus all terms in the identities below are holomorphic at every cusp.

\begin{lemma}[Wronskian--ODE certificate]\label{lem:WH}
For a row of type $R2$, if
\begin{equation}
 K=P_2F^2,\qquad H_2=0,
 \label{eq:WH2}
\end{equation}
then $F(t(q))$ is the normalized analytic solution of $L_2$, and $q$ is its
canonical nome.  For a row of type $R3$, if
\begin{equation}
 K^2=P_3F^2,\qquad H_3=0,
 \label{eq:WH3}
\end{equation}
then $F=u_0^2$, where $u_0$ is the normalized analytic solution of
\eqref{eq:root}, and again $q$ is the canonical nome of $L_3$.
\end{lemma}

\begin{proof}
For $R2$, write $A=tP_2'$ and $R=t(-b+ct)$.  Substituting
$\tht=K^{-1}D$ into $L_2F$ and using $K=P_2F^2$ gives the exact identity
\[
 K^3L_2F=P_2^2F\,H_2.
\]
Thus $L_2F=0$.  Moreover $K=P_2F^2$ is equivalent to
\[
 \tht\log q=\frac1{P_2F^2},
\]
which is precisely the normalized reduction-of-order formula for the ratio of
the logarithmic Frobenius solution to $F$.  Since $t=q+O(q^2)$, the additive
constant is fixed and $q$ is the canonical nome.

For $R3$, put $u=F^{1/2}=1+O(q)$ and use $K^2=P_3F^2$, taking the branch
$K=F\sqrt{P_3}$ at $q=0$.  A direct substitution in \eqref{eq:root} yields
\[
 \frac{\widetilde L u}{u}=\frac{H_3}{4F^4}.
\]
Hence $u$ solves the root equation.  The identity
$\tht\log q=(F\sqrt{P_3})^{-1}$ is its normalized reduction-of-order formula,
so $q$ is canonical.  The symmetric-square relation gives $F=u^2$ as the
analytic solution of $L_3$.
\end{proof}

The identities in Lemma~\ref{lem:WH}, and then the source identity, are now
finite Sturm checks.  The parameters are
\begin{center}
\small
\begin{tabular}{@{}c c c c c@{}}
\toprule
family & type & $(a,b,c)$ & level & $\operatorname{wt}F$\\
\midrule
$\mathbf A$&R2&$(7,2,-8)$&6&1\\
$\mathbf B$&R2&$(9,3,27)$&36&1\\
$\mathbf C$&R2&$(10,3,9)$&6&1\\
$\mathbf E$&R2&$(12,4,32)$&8&1\\
$\mathbf F$&R2&$(17,6,72)$&12&1\\
$\alpha$&R3&$(10,4,64)$&12&2\\
$\gamma$&R3&$(17,5,1)$&6&2\\
$\delta$&R3&$(7,3,81)$&12&2\\
$\varepsilon$&R3&$(12,4,16)$&8&2\\
$\zeta$&R3&$(9,3,-27)$&9&2\\
$\eta$&R3&$(11,5,125)$&20&2\\
\bottomrule
\end{tabular}
\end{center}

\begin{table}[ht]
\centering
\small
\begin{tabular}{@{}c r r r l@{}}
\toprule
family & Wronskian bound & ODE bound & source bound & poles of $t$\\
\midrule
$\mathbf A$&2&6&3&$1/3:-1$\\
$\mathbf B$&12&36&18&$1:-1,\ 1/2:-4,\ 1/4:-1$\\
$\mathbf C$&2&6&3&$1/2:-1$\\
$\mathbf E$&2&6&3&$1/2:-1$\\
$\mathbf F$&4&12&6&$1/2:-2$\\
$\alpha$&8&16&8&$1/2:-2,\ 1/6:-2$\\
$\gamma$&4&8&4&$1/2:-1,\ 1/3:-1$\\
$\delta$&8&16&8&$1:-1,\ 1/2:-2,\ 1/4:-1$\\
$\varepsilon$&4&8&4&$1/2:-1,\ 1/4:-1$\\
$\zeta$&4&8&4&$1/3:-1$\\
$\eta$&12&24&12&$1:-1,\ 1/2:-4,\ 1/4:-1$\\
\bottomrule
\end{tabular}
\caption{Sturm bounds for the three exact identities.  A pole entry
$1/c:-m$ means $\operatorname{ord}_{1/c}t=-m$; in every entry
$\operatorname{ord}_{1/c}F=m$, establishing \eqref{eq:polecancel}.}
\label{tab:bounds}
\end{table}

\begin{proposition}\label{prop:eta-canonical}
For every row of Table~\ref{tab:eta}, the displayed $t,F$ are the canonical
parametrization of the corresponding recurrence.
\end{proposition}

\begin{proof}
Ligozat's criterion gives the stated weights, characters and cusp orders.
The two sides of the Wronskian identity in \eqref{eq:WH2} or
\eqref{eq:WH3} are holomorphic modular forms of weight $2$ or $4$,
respectively.  The forms $H_2,H_3$ are holomorphic of weight $6$ and $8$.
Exact rational Fourier arithmetic verifies the identities through the bounds
in the first two numerical columns of Table~\ref{tab:bounds}.  Sturm's theorem
therefore proves them identically.  Lemma~\ref{lem:WH} finishes the proof.
\end{proof}

\section{The Eisenstein identities}

For a Dirichlet character $\chi$ put
\[
 \sigin{\chi}{k}(m)=\sum_{d\mid m}\chi(m/d)d^k,
 \qquad
 \sigout{\chi}{k}(m)=\sum_{d\mid m}\chi(d)d^k,
\]
and let $\sigma_3(m)=\sum_{d\mid m}d^3$.  A value at a non-integer argument,
such as $\sigma_3(m/2)$, is understood as zero unless the denominator divides
$m$.

\begin{table}[ht]
\centering
\small
\renewcommand{\arraystretch}{1.24}
\begin{tabular}{@{}c p{11.8cm}@{}}
\toprule
family & coefficient $c(m)$ of $E=\sum_{m\ge1}c(m)q^m$\\
\midrule
$\mathbf A$ & $\sigout{\chim}{2}(m)-\sigout{\chim}{2}(m/2)$\\
$\mathbf B$ & $\sigin{\chim}{2}(m)-6\sigin{\chim}{2}(m/2)-8\sigin{\chim}{2}(m/4)$\\
$\mathbf C$ & $\sigin{\chim}{2}(m)-8\sigin{\chim}{2}(m/2)$\\
$\mathbf E$ & $\sigin{\chifour}{2}(m)-8\sigin{\chifour}{2}(m/2)$\\
$\mathbf F$ & $\sigin{\chim}{2}(m)-7\sigin{\chim}{2}(m/2)-8\sigin{\chim}{2}(m/4)$\\[2pt]
$\alpha$ & $\sigma_3(m)-17\sigma_3(m/2)-9\sigma_3(m/3)+16\sigma_3(m/4)+153\sigma_3(m/6)-144\sigma_3(m/12)$\\
$\gamma$ & $\sigma_3(m)-28\sigma_3(m/2)+63\sigma_3(m/3)-36\sigma_3(m/6)$\\
$\delta$ & $\sigma_3(m)-14\sigma_3(m/2)-\sigma_3(m/3)+16\sigma_3(m/4)+14\sigma_3(m/6)-16\sigma_3(m/12)$\\
$\varepsilon$ & $\sigma_3(m)-21\sigma_3(m/2)+84\sigma_3(m/4)-64\sigma_3(m/8)$\\
$\zeta$ & $\displaystyle\sum_{d\mid m}\chim(d)\chim(m/d)d^3$\\
$\eta$ & $\sigin{\chifive}{3}(m)-14\sigin{\chifive}{3}(m/2)-16\sigin{\chifive}{3}(m/4)$\\
\bottomrule
\end{tabular}
\caption{The eleven eta-family Eisenstein sources.}
\label{tab:sources}
\end{table}

Each row is a standard Eisenstein series or a linear combination of its oldform
embeddings.  In the ordinary $E_4$ rows the coefficients of the oldforms sum
to zero, so the constant term cancels.  In the character rows the constant
term is either zero by primitivity or cancels between the displayed oldforms.
Thus Table~\ref{tab:sources} defines a holomorphic form in the same
$M_{r+1}(\Gzero(N),\chi_F)$ as $F\thq t$.

\begin{proposition}\label{prop:eta-source}
For all eleven rows of Table~\ref{tab:eta},
\[
 \boxed{\ F\thq t=E\ },
\]
where $E$ is the Eisenstein series of Table~\ref{tab:sources}.
\end{proposition}

\begin{proof}
At a pole of $t$, equation \eqref{eq:polecancel} shows that $tF$ has order
zero, while $K=D\log t$ is holomorphic.  Hence $F\thq t=tFK$ is holomorphic at
every cusp; it is also holomorphic in the upper half-plane because eta has no
zeros there.  Its weight is $r+1$ and its character is $\chi_F$.
Exact Fourier arithmetic gives $F\thq t-E=0$ through the source Sturm bound in
Table~\ref{tab:bounds}; Sturm's theorem proves the identity.
\end{proof}

\section{The $\mathbf D$ family on $\Gamma_1(5)$}

The remaining order-two sporadic family has $(a,b,c)=(11,3,-1)$ and
$P(t)=1-11t-t^2$.  It is the classical Ap\'ery $\zeta(2)$ family.  Zagier's
exact modular parametrization \cite{Zagier} gives
\begin{equation}
 t=q\prod_{n\ge1}(1-q^n)^{5(\frac n5)},\qquad
 P(t)=t\frac{\eta_1^6}{\eta_5^6},\qquad
 F^2=t^{-1}\frac{\eta_5^5}{\eta_1}.
 \label{eq:D-zagier}
\end{equation}
The corresponding Beauville family is the one with group $\Gamma_1(5)$.
By Proposition~\ref{prop:collapse}, $\Phi=F\thq t=tPF^3$, and
\eqref{eq:D-zagier} therefore gives
\begin{equation}
 \boxed{\ \Phi^2=t\,\eta_1^9\eta_5^3.\ }
 \label{eq:D-square}
\end{equation}

Let $\chi$ be the primitive quartic character modulo $5$, normalized by
$\chi(2)=i$.  Put $\psi_1=\Re\chi$, $\psi_2=\Im\chi$ and
\begin{equation}
 E_D=\sum_{m\ge1}\left(\sum_{d\mid m}(\psi_1(d)-2\psi_2(d))d^2\right)q^m.
 \label{eq:ED}
\end{equation}
For clarity, let $\mathcal E_\chi$ denote the standard weight-three
Eisenstein series whose positive Fourier coefficients are
$\sum_{d\mid m}\chi(d)d^2$.  Then
\[
 E_D=\left(\frac12+i\right)\mathcal E_\chi
     +\left(\frac12-i\right)\mathcal E_{\bar\chi}.
\]
Both summands are holomorphic modular forms on $\Gamma_1(5)$.  With the
standard normalization $B_{3,\chi}=(12+6i)/5$, their constant terms cancel
exactly, so \eqref{eq:ED} is a holomorphic weight-three Eisenstein form with
zero constant term.

\begin{lemma}\label{lem:D-holo}
The right-hand side of \eqref{eq:D-square} is a holomorphic weight-six form
on $\Gamma_1(5)$.
\end{lemma}

\begin{proof}
The eta product $G=\eta_1^9\eta_5^3$ is a weight-six form on
$\Gamma_0(5)$ with quadratic character $\chi_5$, which becomes trivial on
$\Gamma_1(5)$.  Its cusp orders on $X_0(5)$ are $2$ and $1$.
The modular function $h=\eta_1^6/\eta_5^6$ has corresponding orders $1$ and
$-1$.  If a cusp of $X_1(5)$ above the pole cusp of $h$ has ramification
index $e$, then \eqref{eq:D-zagier} and $P(t)\sim-t^2$ force
$\operatorname{ord}t=-e$, whereas $\operatorname{ord}G=e$; the pole cancels.
Above the other cusp, where $h$ has a zero, the same order comparison shows
that $t$ cannot have a pole.  There are no poles in the upper half-plane,
since $h$ is nonzero there and $P(t)=th$ would otherwise have incompatible
pole orders.  Thus $tG$ is holomorphic.
\end{proof}

\begin{proposition}\label{prop:D-source}
For the $\mathbf D$ family, $\Phi=E_D$.
\end{proposition}

\begin{proof}
Both $\Phi^2=t\eta_1^9\eta_5^3$ and $E_D^2$ lie in
$M_6(\Gamma_1(5))$.  A safe Sturm bound is
\[
 \frac6{12}[\mathrm{SL}_2(\Z):\Gamma_1(5)]
 =\frac6{12}\cdot24=12.
\]
The exact expansion
\[
E_D=q-7q^2+19q^3-23q^4+q^5+47q^6-97q^7+105q^8
     -62q^9-7q^{10}+122q^{11}-257q^{12}+O(q^{13})
\]
gives
\begin{align*}
E_D^2={}&q^2-14q^3+87q^4-312q^5+685q^6-794q^7-285q^8\\
&+3308q^9-7441q^{10}+9400q^{11}-4829q^{12}+O(q^{13}).
\end{align*}
Direct expansion of the product on the right of \eqref{eq:D-square} gives
exactly the same coefficients.  Sturm's theorem yields
$\Phi^2=E_D^2$.  Since both $\Phi$ and $E_D$ begin with $q$, the sign is
positive: $\Phi=E_D$.
\end{proof}

\section{Main theorem and scope}

\begin{theorem}[Eisenstein-source theorem]\label{thm:main}
For each of the twelve modularly identifiable sporadic Ap\'ery-like families
\[
 \mathbf A,\mathbf B,\mathbf C,\mathbf D,\mathbf E,\mathbf F,
 \alpha,\gamma,\delta,\varepsilon,\zeta,\eta,
\]
the normalized companion source is
\[
 \Phi=F\thq t,
\]
and is exactly the pure Eisenstein series specified in
Table~\ref{tab:sources}, with the $\mathbf D$ row given by
\eqref{eq:ED}.  In particular, the cuspidal projection of every one of these
twelve sources is zero.
\end{theorem}

\begin{proof}
Proposition~\ref{prop:collapse} identifies the canonical source.
Propositions~\ref{prop:eta-canonical} and \ref{prop:eta-source} prove the
eleven eta cases.  Proposition~\ref{prop:D-source} proves $\mathbf D$.
\end{proof}

\begin{remark}[Literature and novelty scope]\label{rem:novelty}
The second-order side overlaps substantially with Zagier's period analysis:
for the six sporadic $R2$ recurrences, weight-three Eisenstein series and their
Eichler integrals already occur explicitly in his modular treatment
\cite{Zagier}.  We therefore do not claim the discovery of Eisenstein
behaviour for those individual families.  The contribution here is the
uniform identification of the \emph{canonical companion source}
$\Phi=F\theta_qt$, the exact twelve-family source table in one normalization,
and a single finite certification that also proves canonicality of the eta
parametrizations.  Case-by-case priority beyond the cited classical examples
should be assessed separately.
\end{remark}

The theorem does not assert anything about the three additional Cooper rows
$s_7,s_{10},s_{18}$ in their canonical coordinates.  Their canonical nomes
are integral, but the eta parametrizations used above are not available in
that normalization.  Nor does the theorem by itself evaluate every fold
connection value: it identifies the modular source exactly.  The subsequent
passage from an Eisenstein source to a zeta or Dirichlet $L$-value is an
Eichler-integral/connection problem and should be kept logically separate.

\section{Exact certificate}

All finite checks in Tables~\ref{tab:bounds} and \ref{tab:sources} were made
with exact rational arithmetic.  The supplementary script
\texttt{eisenstein\_source\_certificate\_v2.py} constructs the eta products
from finite Euler products, computes $K$ from \eqref{eq:Keta}, checks the
Ligozat congruences and cusp orders, and verifies the Wronskian, ODE and source
identities through their stated Sturm bounds.  It also performs the
$\Gamma_1(5)$ square check through $q^{12}$.  As an implementation guard, the
source identities are additionally checked through $q^{60}$; those extra
coefficients are not used in the proofs.

\begin{thebibliography}{9}
\bibitem{AZ}
G.~Almkvist and W.~Zudilin,
\emph{Differential equations, mirror maps and zeta values},
in Mirror Symmetry V, AMS/IP Stud. Adv. Math. 38 (2006), 481--515;
arXiv:math/0402386.

\bibitem{Beukers}
F.~Beukers,
\emph{Irrationality proofs using modular forms},
Ast\'erisque 147--148 (1987), 271--283.

\bibitem{Ligozat}
G.~Ligozat,
\emph{Courbes modulaires de genre 1},
Bull. Soc. Math. France, M\'emoire 43 (1975).

\bibitem{Sturm}
J.~Sturm,
\emph{On the congruence of modular forms},
Lecture Notes in Math. 1240, Springer (1987), 275--280.

\bibitem{Zagier}
D.~Zagier,
\emph{Integral solutions of Ap\'ery-like recurrence equations},
in Groups and Symmetries, CRM Proc. Lecture Notes 47 (2009), 349--366.
\end{thebibliography}

\end{document}
